\documentclass[11pt]{article}

\usepackage[margin=1in]{geometry}
\usepackage{graphicx}
\usepackage{booktabs}
\usepackage{amsmath}
\usepackage{amssymb}
\usepackage{siunitx}
\usepackage{hyperref}
\usepackage{caption}
\usepackage{subcaption}

\title{Active Homeostatic Compensation in Transformer Grokking:\\Evidence for Le Chatelier's Principle in Neural Networks}
\author{Mat Thompson\\Independent Researcher, Raleigh, NC}
\date{\today}

\begin{document}

\maketitle

\begin{abstract}
Layer 0 suppressors implement factuality hedging in large language models by downweighting confident continuations and boosting hedging tokens. Prior work showed that ablating these heads improves factual preference and calibration, and that structured noise accelerates circuit formation only in sub saturated regimes. What remains unexplained is why the system keeps suppressors near a natural strength and why perturbations in either direction produce a non linear U shaped effect on grokking speed. I test the hypothesis that neural networks maintain a homeostatic equilibrium through active variance redistribution across attention heads. On a parity grokking task I scale a layer 0 suppressor with an omega parameter and then run a kill test in which I either allow or block compensation by freezing the remaining heads at initialization. Across seeds frozen heads either slow learning by up to 4.9 times or drive the network to brittle solutions that grok on the training distribution but fail to generalize. This behaviour is consistent with Le Chatelier's principle applied to neural networks and suggests that suppressor compensation shapes not only how fast circuits form but which circuits crystallize.
\end{abstract}

\paragraph{Contributions.}
\begin{itemize}
  \item \textbf{Kill test for homeostatic compensation.} I design a simple kill test on a parity grokking task and show that blocking compensation in layer 0 either causes severe slowdown or produces brittle non generalizing solutions across seeds.
  \item \textbf{Speed versus stability trade off.} I demonstrate that compensation controls both how fast the network escapes the memorization plateau and which attractor it falls into, with one seed revealing a fast but unstable frozen solution versus a slower but robust free solution.
  \item \textbf{Le Chatelier view of learning.} I frame suppressor compensation as a neural analogue of Le Chatelier's principle and connect the U shaped omega curve to a metastable equilibrium in suppressor strength.
  \item \textbf{Developmental control at layer 0.} I extend the suppressor series by treating layer 0 compensation as a developmental control mechanism that shapes critical windows for circuit formation.
  \item \textbf{Bridge to ironic negation circuits.} I place these results alongside recent work on ironic rebound in transformer models and argue that both phenomena reflect dynamic negotiations between suppression and amplification in different parts of the network.
\end{itemize}

\section{Introduction}

Layer 0 suppressors sit at the first bottleneck in transformer models and regulate how much the network hedges when it is uncertain. Paper 1 showed that a small coalition of early heads systematically downweight factual continuations and boost hedging or editorial tokens \cite{thompson2025suppressors}. Ablating these heads improves factual preference and calibration which ties suppressor strength directly to hallucination behaviour. Paper 2 showed that structured noise is developmental fuel for circuit formation but only in sub saturated regimes and identified observable dependent saturation boundaries in the learning landscape \cite{thompson2025timescale}. In saturated regimes no amount of noise helps. Together these results described the anatomy and dynamics of a single phenomenon: how intelligence crystallizes from a noisy early phase under constraint.

Work on ironic rebound in transformer models lives one layer up the stack. Mann et al.\ show that when a model receives a negation instruction such as ``do not mention X'' it often becomes more likely to produce X than in a neutral baseline and that this effect strengthens with longer or more semantic distractors \cite{mann2025dontthink}. Their circuit analysis finds early heads that suppress forbidden tokens and sparse middle layer heads that amplify them under load. This reveals a dynamic negotiation between suppression and rebound in the middle of the network. In this paper I focus on the homeostatic story at layer 0. Instead of asking how models fail to obey negation under context load I ask how suppressor equilibria shape which circuits crystallize during grokking.

These papers left a central question open. Why is the system's preference for suppressors so strong. In parity grokking experiments I introduced an omega parameter that scales the output of a chosen layer 0 head. The natural setting $\omega = 1.0$ produced the slowest grokking. Perturbing omega in either direction to $0.5$ or $1.5$ accelerated grokking by roughly forty percent. The relationship between suppressor strength and grokking speed was U shaped rather than linear. If weakening suppressors helps and strengthening them also helps, why does the system settle on an intermediate value. Why does it resist learning most strongly at its natural configuration.

The shape of the curve suggests that the system is not merely following different paths through the same loss landscape. It points instead to an equilibrium. At $\omega = 1.0$ the network behaves as if it were in a metastable state. Training pushes on the weights but the system sits in a memorizing configuration for thousands of steps before collapsing into a generalizing circuit. Perturbing omega pushes the system away from this balance and induces instability. Instability accelerates escape from the memorization plateau.

This behaviour is reminiscent of Le Chatelier's principle. When a physical system at equilibrium is perturbed it responds by shifting internal variables in a way that opposes the perturbation. Heat a gas and it expands to reduce the change in temperature. Compress a gas and it develops pressure. The system works to restore balance \cite{kringelbach2024thermodynamics}. If neural networks obey an analogous principle then when I weaken a suppressor head other heads should adjust their outputs to preserve some quantity such as total information flow or attention budget. If I block that adjustment by freezing the other heads in place learning should be severely impaired.

The central hypothesis of this paper is that transformer networks maintain a homeostatic equilibrium through active variance redistribution across attention heads. When I perturb a suppressor by scaling its output the remaining heads adjust to oppose the change. This redistribution is not a passive by product of optimization. It is a coordinated response that preserves a preferred equilibrium state and shapes how the model traverses weight space during training.

To test this hypothesis I designed a kill test on a small transformer trained to grok the parity function. In the baseline condition $\omega = 1.0$ and all heads are free. In the perturbed condition $\omega = 0.5$ and all heads are free so other heads can compensate. In the frozen condition $\omega = 0.5$ and heads 1 through 3 in layer 0 are frozen at their initial values. Head 0 is perturbed and the remaining heads cannot compensate. If compensation is passive drift the frozen runs should grok at similar times to the perturbed runs and differ only in the path taken. If compensation is active and necessary the frozen runs should show severe impairment.

Across five seeds I find that blocking compensation always harms learning. In the cleanest seed frozen heads grok at $10{,}800$ steps compared to $2{,}200$ steps in the perturbed condition, a $4.9$ times slowdown. In other seeds the slowdown is smaller but consistently present. In one seed the frozen run groks fastest but converges to a brittle solution with test accuracy near $0.5$ while the free run groks more slowly but reaches $0.98$ test accuracy. This shows that compensation does not just control how fast the model learns. It shapes which circuits crystallize and whether they generalize beyond the training set.

These results bridge mechanistic interpretability and thermodynamic reasoning. Mechanistically I can identify which heads compensate when suppressors are perturbed and measure how freezing them alters grokking trajectories. Thermodynamically the pattern is consistent with a Le Chatelier style response: when I perturb suppressors the system shifts internal degrees of freedom to restore a preferred equilibrium. When I block those shifts the system either slows dramatically or learns the wrong thing. Understanding neural network learning therefore requires thinking in terms of homeostasis and equilibrium rather than only in terms of gradients and optima.

\section{Methods}

\subsection{Task}

I use a binary parity task as the test bed for the kill test. Each example is a bit string of length $10$ to $12$ sampled uniformly from all possible strings in that range. The label is $1$ if the number of ones in the string is odd and $0$ if it is even. Parity is a Zone 2 task. The model must maintain a running parity state across positions and cannot solve the task by memorizing a static mapping from inputs to outputs. This makes it an ideal setting for studying grokking and circuit formation.

The training set contains $1{,}000$ examples and the test set contains $500$ examples. These sizes are large enough that memorization alone is insufficient but small enough that a small transformer can eventually grok the task under high regularization.

\subsection{Model}

The model is a two layer transformer with four attention heads per layer and a model dimension of $64$. I use learned token and positional embeddings, a standard multi head self attention block, and a two layer feedforward network in each block. The classifier reads from the final token and outputs logits over the two parity classes. There is no dropout. I use AdamW with learning rate $1\times 10^{-3}$ and weight decay equal to $1.0$ to enforce a strong bias toward simpler circuits. All runs use the same architecture and hyperparameters; only the omega scaling and head freezing differ.

\subsection{Omega scaling}

I designate one head in layer 0 as the suppressor head. Its value projection output is multiplied by a scalar $\omega$ before the outputs of all heads are concatenated and passed through the output projection. Omega equal to $1.0$ leaves the head unmodified. Omega equal to $0.5$ halves its contribution. This is the same control parameter used in prior work on omega sweeps which showed that grokking time as a function of omega follows a U shaped curve with the slowest learning at $\omega = 1.0$ and faster learning at both $\omega = 0.5$ and $\omega = 1.5$.

\subsection{Kill test conditions}

For each random seed I train three conditions.

\textbf{Baseline} uses $\omega = 1.0$ with all heads trainable. This is the natural equilibrium configuration.

\textbf{Perturbed} uses $\omega = 0.5$ applied to the suppressor head in layer 0 while all heads remain trainable. Other heads are free to adjust their outputs and compensate.

\textbf{Frozen} uses $\omega = 0.5$ and freezes heads 1, 2, and 3 in layer 0 at their initialization. Gradients into the corresponding slices of the query, key, value, and output projection matrices are zeroed at every step so these heads cannot adapt. Only the perturbed head and other layers remain trainable. This blocks variance redistribution across layer 0 heads and isolates the effect of compensation.

All runs are trained for $50{,}000$ steps with batch size $256$. I log training and test accuracy every $100$ steps and record the first step at which test accuracy exceeds $0.9$ as $T_{\text{grok}}$.

\section{Results}

\subsection{Kill test outcomes across seeds}

Table~\ref{tab:killtest} summarizes $T_{\text{grok}}$ and final test accuracy for the baseline, perturbed, and frozen conditions across four seeds. Seed 4 did not complete due to time constraints and is omitted. In every seed the frozen condition harms learning either by slowing grokking or by driving the network to a brittle non generalizing solution.

\begin{table}[t]
\centering
\begin{tabular}{lllll}
\toprule
Seed & Condition & $T_{\text{grok}}$ & Final test acc & Qualitative pattern \\\midrule
0 & Baseline & 3200 & 0.98 & Clean grokking \\
0 & Perturbed & 2200 & 0.998 & Fast and stable \\
0 & Frozen & 10800 & 0.48 & Strong slowdown \\
1 & Baseline & 2700 & 0.994 & Clean grokking \\
1 & Perturbed & 3400 & 0.998 & Fast and stable \\
1 & Frozen & 3900 & 0.49 & Modest slowdown \\
2 & Baseline & 30000 & 0.74 & Slow baseline \\
2 & Perturbed & 36800 & 0.98 & Slowest but stable \\
2 & Frozen & 3400 & 0.51 & Fast but brittle \\
3 & Baseline & 1400 & 0.96 & Clean grokking \\
3 & Perturbed & 1400 & 0.998 & Clean grokking \\
3 & Frozen & 1700 & 0.998 & Small slowdown \\\bottomrule
\end{tabular}
\caption{Kill test outcomes on parity grokking. $T_{\text{grok}}$ is the first step at which test accuracy exceeds $0.9$. Final test accuracy is measured at step $50{,}000$. Frozen conditions never improve both speed and stability relative to free compensation.}
\label{tab:killtest}
\end{table}

\subsection{Seed zero: strong kill test}

Seed 0 provides the canonical expression of the kill test. The perturbed condition groks at $2{,}200$ steps while the baseline groks at $3{,}200$. Weakening the suppressor accelerates learning in the presence of free compensation. Freezing compensatory heads shifts $T_{\text{grok}}$ out to $10{,}800$ steps, a factor of $4.9$ slower than the perturbed run and $3.4$ times slower than the baseline. The frozen model also finishes with test accuracy near $0.5$, indicating that even after a long search it settles into a weak or unstable generalizing circuit.

This seed validates the original intuition behind the kill test. When suppressor strength is perturbed other heads actively compensate to restore a balanced state. If that compensation is blocked the system must traverse weight space without its usual homeostatic response and learning becomes dramatically less efficient.

\subsection{Kill test learning curves}

Figure~\ref{fig:killtest_curves} shows test accuracy and loss as a function of training step for the three conditions in seed zero. The perturbed run rises smoothly and crosses the grokking threshold early. The baseline run lags behind but still transitions cleanly. The frozen run meanders near chance for many steps before eventually grokking and then exhibits a noisier trajectory. This visualizes the basic effect of compensation. When compensation is allowed the system moves quickly and stably toward the generalizing circuit. When it is blocked the system remains stuck near the memorization plateau and explores a rougher path.

\begin{figure}[t]
\centering
\includegraphics[width=0.9\textwidth]{../../reports/compensation_kill_test.pdf}
\caption{Kill test learning curves for seed zero. Top panel shows test accuracy over training for the baseline, perturbed, and frozen conditions. Bottom panel shows the corresponding loss curves. The dashed horizontal line marks the grokking threshold at test accuracy equal to zero point nine.}
\label{fig:killtest_curves}
\end{figure}

\subsection{Seeds one and three: modest slowdowns}

Seeds 1 and 3 show smaller but consistent slowdowns. In seed 1 the frozen run groks at $3{,}900$ steps compared to $3{,}400$ for the perturbed run and $2{,}700$ for baseline. In seed 3 the frozen run groks at $1{,}700$ steps while baseline and perturbed grok together at $1{,}400$. All three conditions achieve high final test accuracy in these seeds. Blocking compensation therefore does not destroy the ability to find a good circuit, but it does make the path slightly less efficient.

These seeds show that the qualitative effect of freezing compensation generalizes beyond the particular geometry of seed 0. The magnitude varies with initialization but the direction is stable. When compensation is blocked grokking does not become easier.

\subsection{Seed two: fast but brittle frozen solution}

Seed 2 reveals a more subtle phenomenon. Here the frozen run groks first at $3{,}400$ steps but ends with test accuracy near $0.5$. The perturbed run groks much later at $36{,}800$ steps yet reaches $0.98$ test accuracy. The baseline sits in between with $T_{\text{grok}}$ around $30{,}000$ and final accuracy around $0.74$.

In this seed the kill test does not show a pure speed penalty. Instead it separates two different learning trajectories. With compensation blocked the model finds a shortcut that achieves the grokking threshold on the test set early but fails to maintain high accuracy over the remainder of training. The free run follows a slower path but converges to a robust generalizing circuit. This suggests that compensation shapes not only how quickly the system escapes the memorization plateau but which attractor it falls into.

\section{Discussion}

The kill test reveals that compensation plays two distinct roles in transformer grokking. It accelerates learning in favourable regimes and it steers the system toward stable generalizing circuits. When compensation is blocked learning either slows dramatically or converges to brittle solutions that do not generalize.

The strong seed 0 result matches the original picture of Le Chatelier style homeostasis. At $\omega = 0.5$ the perturbed head is weakened and other heads pick up the slack. This redistribution of variance introduces instability into the memorization plateau and accelerates escape. When those compensatory heads are frozen the system can no longer rebalance. It remains stuck near the perturbed configuration and must solve parity by brute force search through weight space. The result is a $4.9$ times slowdown.

Seeds 1 and 3 show that this is not an artefact of a single random initialization. When compensation is blocked the cost can be as small as a twenty percent delay, but in no seed does freezing compensation make learning faster while preserving generalization. The system benefits from being able to redistribute variance across heads even when the magnitude of that benefit is modest.

Seed 2 is more revealing. Here freezing compensatory heads causes the model to grok quickly but badly. $T_{\text{grok}}$ is low and the final test accuracy collapses toward chance. The free run groks much later but reaches high test accuracy. This pattern suggests that homeostatic compensation is not just a speed up mechanism. It constrains the space of solutions the model explores. With compensation active the system is pushed toward circuits that maintain some conserved quantity such as total information flow or representation rank. With compensation blocked the system can satisfy the loss with a narrow shortcut that does not respect that conserved structure.

Viewed through a thermodynamic lens the suppressor heads and their companions define an equilibrium manifold in weight space \cite{kringelbach2024thermodynamics}. Training near this manifold exhibits critical slowing down. Small perturbations decay slowly and the system spends many steps in a metastable memorizing state. Strong perturbations push the system off the manifold. Compensation then acts to pull it back while simultaneously destabilizing the memorization plateau. When compensation is blocked the manifold is broken. The system either crawls along a distorted plateau for many more steps or hops into a different attractor that solves the loss without restoring equilibrium.

\section{Implications}

\subsection{Learning dynamics}

These results flesh out the picture of grokking as a phase transition. The U shaped dependence of grokking time on omega already showed that the natural suppressor configuration is a metastable equilibrium. The kill test shows that compensation is the mechanism that maintains and breaks that equilibrium. With compensation active the system responds to perturbations by redistributing variance and creating the instability required to leave the memorization plateau. With compensation blocked the system either moves more slowly or drops into a qualitatively different solution that fails to generalize.

This suggests that learning curves should be interpreted in light of homeostatic behaviour. Slow learning can indicate that the system is sitting near an equilibrium that resists change. Fast learning can indicate that compensation has destabilized that equilibrium. The seed two result shows that fast grokking is not always a sign of healthy generalization.

The ironic rebound results of Mann et al.\ give a complementary view of the same story in a different regime \cite{mann2025dontthink}. In their setting early layers implement suppression in response to a negation instruction while sparse middle layer heads amplify the forbidden tokens under load. In my setting the key players are layer 0 suppressors and their compensators during grokking. In both cases suppression is not a static property of a single head. It is a dynamic process in which different parts of the circuit pull in opposite directions and the final behaviour depends on how that negotiation resolves over depth and time.

\subsection{Developmental control}

Papers one and two treated suppressors as developmental control knobs. By scaling suppressor strength I could shift the timing of critical windows and acceleration of circuit formation. The present results add a constraint. The system does not passively accept suppressor perturbations. It pushes back. When I weaken suppressors the other heads compensate. When I try to block that compensation the system either slows dramatically or finds brittle shortcuts.

Developmental control therefore cannot rely on one dimensional knobs alone. It must account for the homeostatic response. To delay a dangerous circuit or accelerate a useful one I need to shape the equilibrium manifold itself, not just twist individual parameters. That means designing loss functions and regularization that make the safe circuits the easiest ones to find under the system's natural compensation dynamics.

\subsection{Alignment}

Alignment work often imagines a system that can be constrained by adding rules, penalties, or guard rails. The behaviour observed here suggests that such constraints will be routed around wherever compensation is possible. If a system maintains an internal equilibrium then mechanical constraints that do not change that equilibrium simply divert pressure elsewhere.

A more promising view is to treat alignment as an equilibrium design problem. The question becomes: what quantities does the system try to conserve, and can I construct training regimes in which conserving those quantities implies aligned behaviour. The kill test provides an example at small scale. The suppressor equilibrium appears to favour circuits that generalize on parity. Breaking that equilibrium allows circuits that fit the training data but not the test distribution. Scaling this insight up is an open challenge, but the direction is clear. I need to understand and shape the homeostatic tendencies of my models rather than treating them as static optimizers.

\section{Future work}

Several questions follow directly from these experiments.

First, what quantity is being maintained. The data are consistent with conservation of some measure of information flow or representational capacity across heads, but I have not yet measured these directly. Future work will track mutual information, effective rank, gradient norms, and attention budgets over training to identify which quantity remains stable under compensation.

Second, how general is the effect. Parity is a clean Zone 2 task but narrow. The same kill test should be applied to finite state automata, small algorithmic reasoning tasks, and eventually to layer 0 suppressors in full scale language models. If homeostatic compensation is a general property of multi head attention rather than a quirk of this setup it should appear in those settings as well.

Third, how symmetric is compensation. The present experiments focus on $\omega = 0.5$. Prior omega sweeps show symmetric acceleration at $\omega = 0.5$ and $\omega = 1.5$. Running parallel kill tests at $\omega = 1.5$ and freezing different numbers of heads would reveal whether compensation capacity is symmetric and distributed or asymmetric and concentrated.

Finally, can I harness compensation to steer learning. If compensation both accelerates escape from bad plateaus and favours generalizing circuits I may be able to design schedules that deliberately perturb and then restore suppressor strength to shape developmental windows. Doing so safely will require a clearer picture of the equilibrium manifold and of the failure modes revealed by seeds like seed 2.

\section{Conclusion}

The experiments in this paper support a simple claim. Neural networks maintain homeostatic equilibrium during learning and they do so through active variance redistribution across components. When I perturb a suppressor head other heads compensate. When I block that compensation learning is always harmed, either by slowing dramatically or by converging to brittle non generalizing solutions. This is Le Chatelier's principle in silicon. The system resists perturbation and works to restore a preferred equilibrium, and any intervention has to be understood in that light.


\bibliographystyle{plain}
\bibliography{../references}

\end{document}
